\chapter{Historical compatibility and future simulation}
\label{ch:math}

\section{Define what is observed before inverting it}

Let $D_k$ be the disclosures available under information cutoff $k$. Each
numeric observation $d_i$ has a context: entity or group perimeter, concept,
measurement basis, dimensions, currency, period and precision. These contexts
are part of the observation operator, not additional columns that can be
discarded after extraction.

A candidate history $h$ contains an opening economic and accounting state,
events through the last historical date $T$, contract terms, accounting
policies, consolidation relationships and latent quantities. Let $M$ denote
the declared model scope: its event templates, admissible rule versions and
allowed unknowns. It excludes arbitrary events that the chosen model has no
way to represent. The reporting operator $O_i(h)$ produces the quantity to
which observation $i$ refers, including its basis and perimeter.

Reported numbers are often rounded. If a number $v_i$ is rounded to the nearest
unit $q_i$, an illustrative admissible interval is
$[v_i-q_i/2,v_i+q_i/2)$ under the corresponding tie convention. The actual
rounding convention, scale and XBRL accuracy metadata govern the imported
constraint. A disclosure with unspecified precision does not acquire a
convenient tolerance silently. Write $\mathcal{I}_i$ for the justified
observation interval, and $C_M(h)$ for the joint accounting, contractual and
semantic constraints. Then the historical compatibility set is
\begin{equation}\label{eq:histset}
 \Hist_M(D_k)=\{h\in\mathcal{H}_M:
 C_M(h)\ \land\ O_i(h)\in\mathcal{I}_i\text{ for every applicable }i\}.
\end{equation}
Qualitative assertions enter as reviewed logical conditions, rather than
numeric intervals. Incompatible source contexts are never forced into the
same equation.

Equation~\eqref{eq:histset} distinguishes three outcomes. An empty set means
that the observations and declared model are inconsistent. A singleton means
that this model identifies a history at its chosen resolution. A larger set
preserves unresolved alternatives. An empty set does not prove that the
company's accounts are false: a wrong extraction, missing event type or
incorrect perimeter can cause the inconsistency. A singleton does not prove
that the model contains the real history.

\section{Why published statements cannot identify all transactions}

The simplest derivation uses fixed account mappings and a fixed basis. Let
$z\in\mathbb{R}^n$ contain unknown historical movement amounts. For $m$
disclosed aggregates, let $A\in\mathbb{R}^{m\times n}$ be the aggregation
matrix and $y=Az\in\mathbb{R}^m$. If there is a nonzero vector $v$ with
$Av=0$, then
\[
 A(z+v)=Az+Av=y.
\]
Both movement vectors give the same reports whenever both satisfy the remaining
constraints. Thus the feasible family is the affine set $z+\ker A$ intersected
with those constraints. Accounting requirements can shrink this family, but
do not generally eliminate it. For example, cash collections in a receivables
bridge are constrained by
\[
 AR_T=AR_0+\text{credit sales}-\text{collections}
       -\text{write-offs}+\text{other classified movements}.
\]
If the note does not distinguish acquisitions, currency effects and write-offs,
several collections paths may fit the same closing balance.

There is a second ambiguity beyond aggregation. In
Equation~\eqref{eq:running}, identical historical movements arise from different
contractual reset terms. Even recovering $z$ exactly would leave these terms
unidentified. AccDSL therefore infers economically relevant states and
contract terms alongside journal movements. A linear inverse solver over
statement lines alone cannot solve this second problem.

\section{The accounting state and its invariants}

For one entity book with $n_a$ accounts in one functional currency, let
$b_t\in\mathbb{Q}^{n_a}$ be
signed account balances, with debit amounts positive and credit amounts
negative. Exact decimal amounts can be represented as rationals with a power
of ten denominator. A journal batch $j_t$ contains the signed functional-currency
amounts posted by one atomic event or explicitly grouped event batch. Its
balance condition is
\begin{equation}\label{eq:posting}
 \boldsymbol{1}^{\mathsf T}j_t=0,\qquad b_t=b_{t^-}+j_t.
\end{equation}
If $\boldsymbol{1}^{\mathsf T}b_0=0$, substitution gives
$\boldsymbol{1}^{\mathsf T}b_t=
\boldsymbol{1}^{\mathsf T}b_{t^-}+\boldsymbol{1}^{\mathsf T}j_t=0$.
Induction establishes conservation through every committed batch. Opening
balances must therefore be supplied as a balanced, explained entry or a
balanced imported state with unresolved composition explicitly identified.

This proof establishes a bookkeeping invariant. It does not establish that an
amount belongs in the chosen account, that a contract exists, or that an
accounting standard permits a treatment. Posting validation additionally
checks account roles, ownership, currency conversions, applicability of rules
and event preconditions.

The balance vector includes temporary revenue and expense accounts during
the period. Period closing transfers their net balance to equity through
balanced entries. A reporting function can include current profit in equity
before closing, or use the closed book afterwards, but must not include it
twice. Income and cash-flow reports use classified movements and exclude
mechanical closing transfers from operating flows. Consequently statements
are functions of balances, journals, contract state and reporting rules:
\begin{equation}\label{eq:report}
 Y_{[u,v)}=R(b_{v^-},J_{[u,v)},X_{v^-},\rho,[u,v)).
\end{equation}
Here $\rho$ specifies the reporting basis and perimeter, and $X$ contains
economic state beyond ledger balances. Cash-flow classification is part of
$R$, not recoverable from the closing cash balance alone.

Each book has its own functional currency. Quantities such as shares, tonnes
and foreign-currency units are typed separately and never added to this
balance equation. A foreign-currency posting stores its original units and
the functional-currency valuation with its rate source. Translation into a
group's presentation currency and consolidation adjustments occur in a
separate book under explicit rules.

\section{Simulate events from economic state}

For a chosen history $h$, the state at $T$ is $X_T(h)$. It contains balances,
open contracts, lots or cohorts, unpaid accruals, commitments and other
quantities needed for future evolution. Let $s$ identify a scenario definition,
$U_t^s$ its supplied economic drivers, $a_t$ management's actions, and
$\varepsilon_t$ future innovations. A behavioural policy $\pi$ chooses
$a_t$ using the permitted information $\mathcal{F}_t$:
\begin{align}
 a_t&=\pi(X_t,U_{\leq t}^s,\mathcal{F}_t;\theta),\label{eq:actions}\\
 e_{t+1}&=\mathcal G(X_t,a_t,U_{t:t+1}^s,\varepsilon_{t+1};\theta),\label{eq:events}\\
 (j_{t+1},m_{t+1})&=\mathcal A(e_{t+1},X_t,r_{t+1}),\label{eq:account}\\
 X_{t+1}&=\mathcal T(X_t,e_{t+1},j_{t+1},m_{t+1}).\label{eq:transition}
\end{align}
The generator $\mathcal G$ emits events, such as a sale, rate reset, payment or default.
The accounting function $\mathcal A$ applies rule package $r$ and returns journal
postings plus measurement annotations $m$. The transition $\mathcal T$ updates economic
and ledger state atomically. A default can therefore change the loan's legal
status, expected recoveries, recognised allowance and cash collections through
related, but distinct, records.

This ordering is a reference abstraction. Within a date, event dependencies
and phases specify the actual order. Expectations used to measure an allowance
or provision must be calculated from the information permitted at that date.
They cannot use the realised future path of the simulation. An inner
expectation model may use quadrature, analytic formulas or nested simulation;
its approximation and seed belong in the measurement annotation.

Some scenarios specify a deterministic macroeconomic path while leaving
idiosyncratic defaults random. Others supply a joint stochastic driver model.
A risk-neutral measure used to value a derivative is distinct from the
real-world measure used to simulate business outcomes. Every stochastic
adapter declares its measure and conditioning information. The accounting
function receives the resulting valuation with that declaration.

\section{Preserve sets and distributions separately}

Let $K_{h,s}$ be the probability law of future reports produced by the forward
model conditional on history $h$, scenario $s$ and its declared behavioural
assumptions. A deterministic simulation is the special case of a point mass.
When no probabilities over historical explanations are justified, the natural
answer is a family of laws,
\begin{equation}\label{eq:family}
 \mathcal{K}_{M,s}(D_k)=\{K_{h,s}:h\in\Hist_M(D_k)\}.
\end{equation}
The system can report each explanation, bounds on a deterministic outcome,
or lower and upper event probabilities across this family. Those are different
result types. A range of conditional means is not a confidence interval, and
an envelope of conditional quantiles is not itself a single predictive
distribution.

If a defensible probability measure $\mu(dh\mid D_k,M)$ supported on
$\Hist_M(D_k)$ is supplied, histories and report outcomes must have specified
measurable spaces. Assume that $K_{h,s}$ is a probability kernel: it is a
probability measure on report outcomes for each $h$, and $h\mapsto K_{h,s}(B)$
is measurable for each measurable event $B$. The predictive law is then
\begin{equation}\label{eq:mixture}
 \Prob(Y^s\in B\mid D_k,M)=
 \int K_{h,s}(B)\,\mu(dh\mid D_k,M).
\end{equation}
For discrete explanations this is the usual weighted sum. For continuous
unknowns, integration uses the declared measure over those unknowns, not a
count of solver solutions. Bayesian weights require a prior and a likelihood;
expert weights remain labelled as elicited. A likelihood must not treat a
statement total and all its components as independent repeated evidence.
Deterministic equality constraints require a reduced parameterisation or
proper conditional-measure treatment, rather than an undocumented Gaussian
penalty with near-zero variance.

Suppose the fixed-loan explanation receives weight $w$ and the floating-loan
explanation receives $1-w$, solely as an illustration. Under the rate shock,
\[
 \Prob(I=5)=w,\qquad\Prob(I=7)=1-w,
 \qquad\E[I]=7-2w.
\]
The expected expense is useful but hides the two different outcomes. For a
general scalar output with finite second moment, put $m(h)=\E[Y\mid h]$.
Write $Y-\E[Y]=(Y-m(h))+(m(h)-\E[Y])$. The cross term vanishes because
$\E[Y-m(h)\mid h]=0$. Taking expectations of the squared expression gives
\begin{equation}\label{eq:variance}
 \operatorname{Var}(Y)=
 \E_h[\operatorname{Var}(Y\mid h)]+
 \operatorname{Var}_h(\E[Y\mid h]).
\end{equation}
The first term is uncertainty within a chosen explanation; the second captures
differences between explanations. This decomposition requires the measure
$\mu$. Without it, the system reports the family in
Equation~\eqref{eq:family}, rather than manufacturing a variance.

Weights themselves can be uncertain. For an admissible class $\mathcal{W}$
of measures, AccDSL can compute
\[
 \underline p(B)=\inf_{\mu\in\mathcal{W}}\int K_{h,s}(B)\,d\mu(h),
 \qquad
 \overline p(B)=\sup_{\mu\in\mathcal{W}}\int K_{h,s}(B)\,d\mu(h).
\]
If every measure supported on the feasible histories is allowed and point
masses are included, these bounds reduce to the infimum and supremum of
$K_{h,s}(B)$ over histories: averages cannot lie outside the range of their
integrands, and point masses attain, or approximate, the extremes. Narrower
weight restrictions require an explicit justification and a separate
optimisation problem.

\section{Resolution and computational limits}

For a requested horizon $H$ and scenario family $\mathcal{S}$, histories may be
grouped only when they have the same conditional law for all requested
outputs:
\[
 h_1\sim_{H,\mathcal{S}}h_2
 \quad\Longleftrightarrow\quad
 K_{h_1,s}=K_{h_2,s}\text{ on those outputs for every }s\in\mathcal{S}.
\]
This definition explains why individual sales can sometimes be replaced by
cohorts, while fixed and floating loans usually cannot be merged for an
interest-rate stress. Exact equivalence can be difficult to prove. A grouping
justified only by approximate numerical agreement carries an error criterion
and is not labelled exact. Source provenance remains available even when
execution shares a common state.

Bounds require comparable honesty. A linear program can provide certified
extrema for a supported linear model. A set of feasible sampled histories
provides observed extrema: its minimum is no smaller than the true infimum,
and its maximum is no larger than the true supremum. It therefore supplies
an inner exploration of possibilities, not a certified outer bound. Interval
methods may provide outer bounds that no single history attains. Every result
records which of these meanings applies.

The initial implementation offers exact bounded choices and linear constraints,
symbolic continuous families such as $\alpha\in[0,1]$, and explicit scenario
execution at selected feasible states. General nonlinear partial identification,
automatic model discovery and exhaustive enumeration of all economic histories
are research extensions, not promises made by the core language.
